\documentclass[11pt, a4paper]{article}

% ── Packages ──────────────────────────────────────────────────────────
\usepackage[utf8]{inputenc}
\usepackage[T1]{fontenc}
\usepackage{lmodern}
\usepackage[margin=1in]{geometry}
\usepackage{amsmath, amssymb, amsthm}
\usepackage{booktabs}
\usepackage{longtable}
\usepackage{graphicx}
\usepackage{hyperref}
\usepackage{xcolor}
\usepackage{listings}
\usepackage{tcolorbox}
\usepackage{polya}

\hypersetup{
  colorlinks=true,
  linkcolor=polyablue,
  urlcolor=polyablue,
  citecolor=polyablue,
}

% ── Code listing style ────────────────────────────────────────────────
\lstdefinestyle{polya-python}{
  language=Python,
  basicstyle=\ttfamily\small,
  keywordstyle=\color{polyablue}\bfseries,
  commentstyle=\color{polyagray}\itshape,
  stringstyle=\color{polyagreen},
  numbers=left,
  numberstyle=\tiny\color{polyagray},
  numbersep=8pt,
  frame=single,
  framerule=0.4pt,
  rulecolor=\color{polyagray},
  breaklines=true,
  showstringspaces=false,
  tabsize=4,
}
\lstset{style=polya-python}
\title{Customer Survival -- SaaS Subscription Churn}
\author{uber-polya}
\date{2026-02-26}

\begin{document}

\maketitle
\tableofcontents
\newpage

% ── Phase A: Problem Formulation ─────────────────────────────────────
\section{Problem Formulation}

\begin{polyamodel}

\textbf{Problem Type}: Find \hfill
\textbf{Domain}: Statistical Inference


\end{polyamodel}

% ── Universe ──────────────────────────────────────────────────────────
\subsection{Universe}

\begin{itemize}
  \item $Kaplan Meier: \{median_basic, median_pro\}$
  \item $Logrank Test: \{statistic, p_value\}$
  \item $Ph Assumption: \{plan_schoenfeld_p, tickets_schoenfeld_p, holds\}$
\end{itemize}

% ── Variables ─────────────────────────────────────────────────────────
\subsection{Variables}

\begin{longtable}{lll}
\toprule
\textbf{Symbol} & \textbf{Meaning} & \textbf{Type / Range} \\
\midrule
\endhead
$\mu_1, \mu_2$ & population means for groups 1 and 2 & $\mathbb{{R}}$ \\
$\bar{{x}}_1, \bar{{x}}_2$ & sample means & $\mathbb{{R}}$ \\
$s_1, s_2$ & sample standard deviations & $\mathbb{{R}}_{{>0}}$ \\
$n_1, n_2$ & sample sizes & $\mathbb{{Z}}^+$ \\
$t$ & test statistic (hypothesis test) & $\mathbb{{R}}$ \\
$p$ & p-value & $[0, 1]$ \\
$\alpha$ & significance level & $[0, 1]$ \\
\bottomrule
\end{longtable}

% ── Structure ─────────────────────────────────────────────────────────
\subsection{Structure}

- **Customers**: 500 SaaS subscribers with subscription duration (months) and churn status
- **Censoring**: \textasciitilde{}30\% of customers are right-censored (still active, have not churned)
- **Plan types**: Basic (n=300, median survival \textasciitilde{}12 months) and Pro (n=200, median survival \textasciitilde{}24 months)
- **Covariate**: Support ticket count (higher tickets correlate with higher churn risk)
- **Questions**: Do Basic and Pro plans have different churn rates? What factors predict churn?

% ── Mapping ───────────────────────────────────────────────────────────
\subsection{Real-World Mapping}

\begin{longtable}{ll}
\toprule
\textbf{Real-World Concept} & \textbf{Mathematical Object} \\
\midrule
\endhead
group average & $sample mean x-bar$ \\
group variability & $sample std s$ \\
sample count & $n$ \\
significance threshold & $alpha = 0.05$ \\
\bottomrule
\end{longtable}

% ── Constraints ───────────────────────────────────────────────────────
\subsection{Constraints}

\begin{enumerate}
  \item $H_0: \mu_1 = \mu_2$
  \hfill \textit{(null hypothesis)}
  \item $H_1: \mu_1 \neq \mu_2$
  \hfill \textit{(alternative hypothesis)}
  \item $\alpha = 0.05$
  \hfill \textit{(significance level)}
  \item $t = \frac{\bar{x}_1 - \bar{x}_2}{\sqrt{\frac{s_1^2}{n_1} + \frac{s_2^2}{n_2}}}$
  \hfill \textit{(Welch's t-test statistic)}
\end{enumerate}

% ── Objective / Claim ─────────────────────────────────────────────────
\subsection{Objective}

\begin{equation}
\text{Reject } H_0 \text{ if } p < \alpha
\end{equation}

% ── Approach ──────────────────────────────────────────────────────────
\subsection{Approach}

\begin{description}
  \item[Suggested approach:] Kaplan-Meier + Log-Rank Test + Cox Proportional Hazards
  \item[Complexity class:] 
  \item[Available tools:] Kaplan-Meier + Log-Rank Test + Cox Proportional Hazards
\end{description}
% ── Phase B: Solution ────────────────────────────────────────────────
\section{Solution}

\begin{polyaresult}

\textbf{Answer}: Statistical test complete


\textbf{Optimal}: No / Unknown \hfill
\textbf{Feasible}: Yes
\textbf{Algorithm}: Kaplan-Meier + Log-Rank Test + Cox Proportional Hazards \quad $$

\textbf{Time}: 0.09642017504666s

\textbf{Certificate}: Assumptions checked; test statistic independently verified.

\end{polyaresult}

% ── Solution Details ──────────────────────────────────────────────────
\subsection{Solution Details}

Statistical analysis computed.

% ── Verification ──────────────────────────────────────────────────────
\subsection{Verification}

\begin{longtable}{lcl}
\toprule
\textbf{Check} & \textbf{Status} & \textbf{Detail} \\
\midrule
\endhead
Km Basic Monotonic & \passmark & Passed \\
Km Pro Monotonic & \passmark & Passed \\
Km Basic In 01 & \passmark & Passed \\
Km Pro In 01 & \passmark & Passed \\
Logrank Significant & \passmark & Passed \\
Cox Plan Hr Gt 1 & \passmark & Passed \\
Cox Tickets Hr Gt 1 & \passmark & Passed \\
Concordance Gt 055 & \passmark & Passed \\
Median Basic Lt Pro & \passmark & Passed \\
Censoring Rate Approx 30Pct & \passmark & Passed \\
All Passed & \passmark & Passed \\
\bottomrule
\end{longtable}

% ── Phase C: Interpretation ──────────────────────────────────────────
\section{Interpretation}

% ── The Question & The Answer ─────────────────────────────────────────
\subsection{The Question}
- **Customers**: 500 SaaS subscribers with subscription duration (months) and churn status
- **Censoring**: \textasciitilde{}30\% of customers are right-censored (still active, have not churned)
- **Plan types**: Basic (n=300, median survival \textasciitilde{}12 months) and Pro (n=200, median survival \textasciitilde{}24 months)
- **Covariate**: Support ticket count (higher tickets correlate with higher churn risk)
- **Questions**: Do Basic and Pro plans have different churn rates? What factors predict churn?.

\subsection{The Answer}

\begin{polyainsight}[title={Bottom Line}]
Test complete.
\end{polyainsight}

% ── What This Means ───────────────────────────────────────────────────
\subsection{What This Means}

- Kaplan-Meier: Monotonically decreasing survival curves for Basic and Pro plans
- Median survival: Basic \textasciitilde{}12 months, Pro \textasciitilde{}24 months
- Log-rank test: p < 0.05 (significant difference between plans)
- Cox PH hazard ratios: Basic vs Pro HR > 1.0 (Basic has higher hazard), tickets HR > 1.0 (more tickets = higher churn)
- 95\% CI: Hazard ratio confidence intervals excluding 1.0 for both covariates
- Concordance index: > 0.55 (better than random prediction)
- PH assumption: Schoenfeld residuals test for proportional hazards

1. Kaplan-Meier estimator: Non-parametric survival function estimate per plan group, handling right-censored observations
2. Log-rank test: Non-parametric hypothesis test comparing survival distributions between Basic and Pro plans
3. Cox Proportional Hazards: Semi-parametric regression modeling hazard as a function of plan type and support ticket count; reports hazard ratios with 95\% confidence intervals
4. Diagnostics: Proportional hazards assumption check via Schoenfeld residuals; concordance index for predictive accuracy

% ── Sensitivity Analysis ──────────────────────────────────────────────

% ── Figures ───────────────────────────────────────────────────────────

% ── Recommendations ───────────────────────────────────────────────────
\subsection{Recommendations}

\begin{enumerate}
  \item Report effect sizes alongside p-values for practical significance assessment.
\end{enumerate}

% ── Limitations ───────────────────────────────────────────────────────
\subsection{Limitations}

\begin{polyawarnbox}
\begin{itemize}
  \item Results are conditional on model assumptions (normality, independence, etc.).
  \item Statistical significance does not imply practical importance.
\end{itemize}
\end{polyawarnbox}

% ── Appendix ─────────────────────────────────────────────────────────

\end{document}